% !PRINTED 493
\weberpage{493}
\weq{13}{\begin{aligned}
n &= 11,\quad & c &= \infty,\ 0,\ 1,\ 2,\ 3,\ 4,\ 5,\ 6,\ 7,\ \ 8,\ 9,\ 10,\\
&& c' &= \infty,\ 2,\ \ 3,\ 4,\ 5,\ 6,\ 7,\ 8,\ 9,\ 10,\ 0,\ \ 1,\\
&& c'' &= 0,\ \infty,\ 10,\ 5,\ 7,\ 8,\ 2,\ 9,\ 3,\ \ 4,\ 6,\ \ 1,\\
&& c''' &= 1,\ 10,\ \ 0,\ 4,\ 6,\ 5,\ 8,\ 7,\ 9,\ \ \ 2,\ 3,\ \infty.
\end{aligned}}

Als Wurzeln der Resolventen 7ten und 11ten Grades können wir nach \S\,82 je eine der beiden folgenden Funktionen betrachten:

\weq{14}{\begin{aligned}
n &= 7,\\
w_\nu &= \frac{(v_\infty - v_\nu)(v_{\nu+1}-v_{\nu+3})(v_{\nu+2}-v_{\nu+6})(v_{\nu+4}-v_{\nu+5})}{i\sqrt{7}\,u^4},\\
w'_\nu &= \frac{(v_\infty - v_\nu)(v_{\nu+1}-v_{\nu+5})(v_{\nu+2}-v_{\nu+3})(v_{\nu+4}-v_{\nu+6})}{-i\sqrt{7}\,u^4},
\end{aligned}}

\weq{15}{\begin{aligned}
n &= 11,\quad w_\nu =\\
&\frac{(v_\infty-v_\nu)(v_{\nu+1}-v_{\nu+2})(v_{\nu+3}-v_{\nu+6})(v_{\nu+4}-v_{\nu+8})(v_{\nu+5}-v_{\nu+10})(v_{\nu+9}-v_{1+7})}{i\sqrt{11}\,u^6},\\
w'_\nu &=\\
&\frac{(v_\infty-v_\nu)(v_{\nu+1}-v_{\nu+6})(v_{\nu+3}-v_{\nu+7})(v_{\nu+4}-v_{\nu+2})(v_{\nu+5}-v_{\nu+8})(v_{\nu+9}-v_{\nu+10})}{-i\sqrt{11}\,u^6},
\end{aligned}}
% [sic] printed last factor is $(v_{\nu+9}-v_{1+7})$; the index pattern would require $v_{\nu+7}$.

worin $\sqrt{7}$ und $\sqrt{11}$ positiv genommen sein sollen.

Es tritt nun zunächst der folgende bemerkenswerte Unterschied zwischen beiden Fällen hervor.

Durch die beiden Vertauschungen $(c')$, $(c'')$ werden im Falle $n = 7$ die $w_\nu$ nur untereinander vertauscht, und zwar in folgender Weise:

\[
\begin{aligned}
& w_0,\ w_1,\ w_2,\ w_3,\ w_4,\ w_5,\ w_6,\\
(c') \quad & w_2,\ w_3,\ w_4,\ w_5,\ w_6,\ w_0,\ w_1,\\
(c'') \quad & w_0,\ w_2,\ w_1,\ w_6,\ w_4,\ w_5,\ w_3,
\end{aligned}
\]

während [nach (13)] im Falle $n = 11$ durch die Vertauschung $(c')$ die Größen $w$ gleichzeitig ihr Vorzeichen ändern, so daß folgende Vertauschungen eintreten:

\[
\begin{aligned}
& w_0,\ w_1,\ w_2,\ w_3,\ w_4,\ w_5,\ w_6,\ w_7,\ w_8,\ w_9,\ w_{10},\\
(c') \quad & -w_2,\ -w_3,\ -w_4,\ -w_5,\ -w_6,\ -w_7,\ -w_8,\ -w_9,\ -w_{10},\ -w_0,\ -w_1,\\
(c'') \quad & w_0,\ w_9,\ w_6,\ w_3,\ w_5,\ w_4,\ w_2,\ w_8,\ w_7,\ w_1,\ w_{10},
\end{aligned}
\]

und ebenso verhält es sich mit $w'$.

Daraus folgt nun (\S\,54), daß die $w_\nu$ Wurzeln einer Gleichung 7ten oder 11ten Grades:

\weq{16}{w^n + A_1 w^{n-1} + A_2 w^{n-2} + \cdots + A_n = 0}
%CONT>>
